\section{Introduction}

Despite the maturity of KV cache eviction methods like H2O and StreamingLLM, a fundamental question remains unanswered: do all tasks respond equally to compression, or does task structure dictate optimal strategy? This question has profound implications for memory-constrained deployment---choosing the wrong compression approach may cost significant accuracy, while practitioners currently lack any principled basis for selection.

The memory challenge is well understood. KV cache memory grows linearly with context length, consuming hundreds of megabytes for even moderate contexts. Methods like H2O's heavy-hitter eviction and StreamingLLM's attention sink retention provide effective compression, achieving 5x or greater memory reduction with minimal accuracy loss in aggregate benchmarks. Quantization offers an orthogonal dimension, reducing precision from 16-bit to 8-bit or 4-bit.

However, these methods share a critical assumption: \emph{one compression strategy fits all tasks}. H2O applies the same heavy-hitter eviction ratio regardless of whether the task is multi-document question answering---requiring cross-document reasoning---or single-document extraction. StreamingLLM uses a fixed window size whether the task processes sequential narrative or requires precise retrieval. This uniformity may leave significant accuracy on the table.

The gap in existing work is striking. While individual methods are extensively evaluated, no systematic study quantifies whether tasks genuinely respond differently to compression. Do some tasks tolerate aggressive eviction while others collapse? Does quantization affect summarization differently than code completion? Without answers, practitioners resort to trial-and-error or conservative choices.

We address this gap by providing the first rigorous quantification of task-dependent KV cache compression behavior. Our key insight is that LongBench tasks cluster into statistically distinct compression response groups---not through intuition but through formal gap statistic analysis revealing $k^*=3$ clusters with gap exceeding standard error. Furthermore, attention entropy computed from the first 100 tokens discriminates task domains with large effect size ($\eta^2=0.522$), suggesting attention patterns encode task structure that correlates with compression tolerance.

Building on this insight, we make the following contributions:

\begin{enumerate}
    \item \textbf{First quantified task-compression clustering.} We demonstrate that 21 LongBench tasks form 3 distinct clusters in their response to 6 compression configurations, validated by gap statistic with $B=500$ bootstrap samples. The clusters are interpretable: Multi-doc QA + Code, Single-doc QA + Few-shot, and Summarization + Synthetic.

    \item \textbf{Attention entropy as task discriminator.} We show that first-100-token attention entropy varies significantly across task domains (F=38.05, $p=2.92 \times 10^{-26}$, $\eta^2=0.522$), validating entropy as a potential routing signal without requiring task labels.

    \item \textbf{Methodological contribution.} We identify that standard accuracy metrics are inadequate for KV compression evaluation when answer formats are diverse, motivating improved evaluation protocols for future work.

    \item \textbf{Empirical foundation for task-conditioned compression.} Our findings establish that task-aware strategy selection is grounded in measurable structure, enabling principled compression decisions rather than one-size-fits-all defaults.
\end{enumerate}
